\addchap{\lsAcknowledgementTitle} 

Die vorliegende Arbeit stellt eine geringfügig überarbeitete Fassung meiner im März 2024 in Leipzig eingereichten Dissertation dar. Meinen Dank möchte ich zunächst der Universität Leipzig aussprechen, die mir durch einen Doktorandenförderplatz die Möglichkeit gegeben hat, mich mit viel Zeit und Freiheit meinem Forschungsthema zu widmen – keine Selbstverständlichkeit im Rahmen einer Promotionsstelle.

Mein größter Dank gilt meinen Betreuern Oliver Czulo (Leipzig) und Noah Bubenhofer (Zürich). Von meinem Doktorvater Oliver Czulo erhielt ich den Anstoß zu einer vergleichenden Untersuchung des Framings von Extremismusvarianten und konnte auf Vorarbeiten von ihm und seinen Kollegen aufbauen. Den Bereich der korpuslinguistischen Diskursanalyse und der Frame-Semantik habe ich erst als Praktikant an seinem Lehrstuhl kennengelernt. Für seine stete Unterstützung, die vielen hilfreichen Gespräche und auch seine Offenheit dafür, das von ihm übernommene Thema ganz in meinem eigenen Sinne auszuarbeiten, danke ich ihm von Herzen.

Noah Bubenhofer gilt mein herzlicher Dank zunächst dafür, dass er mich gleich zwei Mal im Laufe meiner Promotionsphase zu Gastaufenthalten in Zürich willkommen geheißen hat. Seine Art, sich korpuslinguistisch mit Diskursen auseinander zu setzen, hat meine Arbeit sehr geprägt und mir an vielen Stellen neue Perspektiven und Ideen gegeben. Die Aufenthalte in Zürich waren aber nicht nur aufgrund des fachlichen Austauschs mit ihm und seinem Team bereichernd. Ich habe mich immer sehr willkommen gefühlt und neben den Impulsen zu meiner Arbeit auch neue Freundschaften gewinnen können. Dafür danke ich Daniel, Larissa, Livia, Maaike, Marina und – gerade auch für seine fruchtbaren methodischen Hinweise und geduldigen Erklärungen zu computerlinguistischen Feinheiten – Niclas.

Dem CRETA-Verein und insbesondere Melanie Andresen und Janis Pagel danke ich dafür, dass sie mir in einem Coaching in der frühen Phase der Promotion wertvolle Hinweise gaben, um mir einen Weg durch den Methodendschungel der reflektierten Textanalyse zu bahnen. Meinen neuen Kollegen am IDS, besonders Marc Kupietz, danke ich für ihre Geduld und die Möglichkeit, die letzten Seiten dieser Dissertation in Mannheim schreiben zu können.

Anna, Kais, Linnéa und Lisa haben viele Stunden mit der gründlichen Lektüre meiner Arbeit zugebracht und mich auf diverse Tippfehler und unübersichtliche Sätze hingewiesen. Dafür danke ich ihnen wirklich sehr.

Zuletzt möchte ich all den Menschen aus meinem privaten Umfeld danken, die mich auf meinem Weg immer unterstützt und begleitet haben – selbst wenn sie nicht wussten, wo eine korpuslinguistische Analyse semantischer Frames überhaupt hinführen sollte: meinen Eltern Silvia und Bernd, meinem Großvater Karl-Heinz und meinem Freund Johannes.
